\section{Two competing trees}
\label{sec:competing-trees}

In a rooted tree $T$, the depth of a vertex is its distance from the root;
write $V_\ell(T)$ for the set of vertices at depth $\ell$.  A
\emph{marking} is a set of vertices at positive depths.  It \emph{hits} a
path if the path contains a marked vertex.  The following construction
provides two trees that each admit many hits along every path, but cannot
both be hit when their depth quotas are shared.

Depth-constrained path hitting appears in the tree formulation of resource
minimization for fire containment studied by Chalermsook and
Chuzhoy~\cite[Section~2]{ChalermsookChuzhoy2010}.
The obstruction below uses the same disjoint-branch counting principle as
their construction in Section~6: more disjoint branches than marks
available in their depth range leave one branch unhit.
Here the quotas vary with depth, and retaining paths according to their
red or blue vertex counts produces separate multiple-hit markings but
excludes a shared one-hit marking.
We prove this two-tree property in full.

\begin{lemma}[Competing trees]\label{lem:competing-trees}
For every integer $d\ge1$, there are two finite rooted trees $T^+,T^-$,
an integer $L\ge1$, and positive integer quotas $p_1,\ldots,p_L$ with the
following properties.
\begin{enumerate}[beginpenalty=10000]
\item All leaves of both trees have depth $L$.  Each tree $T^\sigma$,
$\sigma\in\{+,-\}$, has a marking $\mathcal M^\sigma$ such that
\begin{equation}\label{eq:tree-separate-marking}
 \begin{aligned}
 |\mathcal M^\sigma\cap V_\ell(T^\sigma)|&\le p_\ell
 &&(1\le\ell\le L),\\
 |\mathcal M^\sigma\cap\pi|&\ge d
 &&\text{for every root-to-leaf path $\pi$ in $T^\sigma$}.
 \end{aligned}
\end{equation}
\item For any markings $\mathcal A^+$ and $\mathcal A^-$ satisfying the
combined bounds
\begin{equation}\label{eq:tree-combined-quota}
 |\mathcal A^+\cap V_\ell(T^+)|+
 |\mathcal A^-\cap V_\ell(T^-)|\le p_\ell
 \qquad(1\le\ell\le L),
\end{equation}
at least one of the two trees has a root-to-leaf path that is not hit by
its marking.
\end{enumerate}
The trees, quotas, and separate markings are constructible by a finite
deterministic procedure depending only on $d$.
\end{lemma}

\begin{proof}
We first build one auxiliary tree $\mathcal T$, with some vertices colored
red or blue.  We then obtain $T^+$ and $T^-$ by retaining suitable complete
paths in two copies of $\mathcal T$.

\paragraph{Phases and global depths.}
Begin with a single root at depth $h_0=0$, and perform $2d-1$ phases.
At the start of phase $a$, let the current leaves be
$u_1,\ldots,u_N$, all at depth $h_{a-1}$.  The phase will add $N$ levels,
so set $h_a=h_{a-1}+N$.  Give the new global depths the quotas
\begin{equation}\label{eq:tree-phase-quotas}
 p_{h_{a-1}+i}=2^{i-1}=\frac{M_i}{2},
 \qquad M_i=2^i,\qquad 1\le i\le N.
\end{equation}
The indices $1,\ldots,N$ are local to this phase.  The successive depth
blocks $(h_{a-1},h_a]$ are disjoint, so each positive depth receives its
quota exactly once.

For each $i$, attach to $u_i$ exactly $M_i$ paths of $i$ edges, called
\emph{spines}.  They share their starting vertex $u_i$ and are otherwise
vertex-disjoint.  Color half of their endpoints red and half blue.  Each
such endpoint has depth $h_{a-1}+i$.  If $i<N$, attach to each endpoint
exactly
\begin{equation}\label{eq:tree-continuations}
 1+\sum_{j=i+1}^{N}2^{j-1}
\end{equation}
paths of $N-i$ edges, sharing only that endpoint.  Call these paths the
\emph{continuation chains}.  All vertices introduced in distinct spines
or continuation chains are new, except for the explicitly shared starting
vertices.  In particular, constructions below different $u_i$ are
disjoint.  No newly introduced vertex on a continuation chain is colored.
For $i=N$, the spine endpoint already has depth $h_a$ and no continuation
is added.

Every new leaf has depth $h_a$, and every vertex at a smaller depth has
a descendant at depth $h_a$.  This proves the equal-depth invariant for
the next phase.  After the last phase put $L=h_{2d-1}$ and let
$\mathcal T$ be the resulting tree.  All its leaves have depth $L$, and
$L\ge1$.

\paragraph{An obstruction to arbitrary markings.}
Fix any marking $\mathcal A$ of $\mathcal T$ with at most $p_\ell$ marks
at each depth $\ell$.  The marks may be anywhere at positive depths;
there is no restriction to colored vertices.  We construct a path from
the root to depth $L$ that avoids $\mathcal A$.

Suppose that a path avoiding $\mathcal A$ has reached a phase-start leaf
$u_i$.  The total number of marks anywhere in the first $i$ new levels
of this phase is at most
\begin{equation}\label{eq:tree-spine-budget}
 \sum_{j=1}^{i}p_{h_{a-1}+j}
 =\sum_{j=1}^{i}2^{j-1}
 =2^i-1=M_i-1.
\end{equation}
There are $M_i$ spines below $u_i$, disjoint after their unmarked starting
vertex.  Hitting each spine would require at least $M_i$ different marks
in these levels.  Equation~\eqref{eq:tree-spine-budget} therefore leaves
one spine entirely unmarked, including its endpoint.

If $i=N$, this spine already reaches the phase boundary.  If $i<N$, the
number of marks in all the remaining new levels is at most
$\sum_{j=i+1}^{N}2^{j-1}$.  By
Equation~\eqref{eq:tree-continuations}, the unmarked endpoint has one more
continuation chain than this bound.  The chains are disjoint after that
endpoint, so one of them is entirely unmarked.  Following it extends the
unhit path to depth $h_a$.

The root cannot be marked.  Starting there and applying this argument
through all phases gives an unhit root-to-leaf path in $\mathcal T$.
Thus no quota-respecting marking of $\mathcal T$ hits every complete path.
The argument applies to a marking fixed on the entire final tree: each
phase uses only the marks in its own block of global depths.

\paragraph{Color counts and pruning.}
At depth $h_{a-1}+i$, the only colored vertices introduced in phase $a$
are the endpoints of the $M_i$ spines below $u_i$.  Spines below the other
phase-start leaves have their endpoints at different depths, and their
descendants are disjoint from those below $u_i$.  Since the depth blocks
of different phases are disjoint, there are exactly $p_{h_{a-1}+i}$ red
vertices and the same number of blue vertices at this global depth.

Every root-to-leaf path in $\mathcal T$ encounters exactly one colored
endpoint in each phase.  It therefore contains exactly $2d-1$ colored
vertices, at least $d$ of which have the same color.  Let
$\mathcal P^+$ be the family of complete paths containing at least $d$
red vertices, and let $\mathcal P^-$ be the analogous family for blue.
These families cover all complete paths.  They are both nonempty:
at every phase-start leaf, one can choose a spine of either color and
continue to the next phase boundary, so in particular there are paths
choosing red in every phase and paths choosing blue in every phase.

In separate copies of $\mathcal T$, retain the union of the paths in
$\mathcal P^+$ and in $\mathcal P^-$, respectively; these are $T^+$ and
$T^-$.  Retain every vertex and edge on each chosen path, including
vertices with only one retained child.  No edge is contracted.  Hence
all original depths are preserved.  Every retained vertex lies on a
chosen complete path, so no vertex of depth less than $L$ becomes a leaf.
Conversely, every retained root-to-leaf path is the unique path to one
of the chosen original leaves, and thus belongs to its specified family.

Mark all retained red vertices in $T^+$ and all retained blue vertices
in $T^-$.  Pruning can only decrease their numbers at each depth.
The preceding color counts give the quotas in
Equation~\eqref{eq:tree-separate-marking}, and the definition of each
retained path family gives at least $d$ marks along every path.  All these
marks have positive depth.

\paragraph{Projection of a combined marking.}
Suppose, for a contradiction, that markings $\mathcal A^+,\mathcal A^-$
satisfy Equation~\eqref{eq:tree-combined-quota} and hit every complete path
in both trees.  Let $\varphi^\sigma:V(T^\sigma)\to V(\mathcal T)$ be the
map identifying each retained copy vertex with its original vertex, and
set
\begin{equation}\label{eq:tree-projected-marking}
 \mathcal A=
 \varphi^+(\mathcal A^+)\cup\varphi^-(\mathcal A^-).
\end{equation}
These maps preserve depths.  Taking a union merges any marks projected
onto the same original vertex, so for every $1\le\ell\le L$,
\[
 |\mathcal A\cap V_\ell(\mathcal T)|
 \le |\mathcal A^+\cap V_\ell(T^+)|
     +|\mathcal A^-\cap V_\ell(T^-)|
 \le p_\ell.
\]
The roots are unmarked, and thus $\mathcal A$ is a valid marking of
$\mathcal T$.  Every complete path in $\mathcal T$ belongs to at least
one of $\mathcal P^+,\mathcal P^-$.  Its retained copy is hit, so its
original path contains a vertex of $\mathcal A$.  This contradicts the
obstruction proved above, and establishes the combined-quota assertion.

Every phase adds finitely many explicitly specified vertices and edges.
There are finitely many phases, and the final pruning is determined by
counting colors along the finitely many complete paths.  This is a finite
deterministic construction from $d$, as claimed.
\end{proof}
